% Book 3, Year 1, chapter 13 -- asymptotic analysis and recurrences.
% Provenance (never printed): sources/book3/BRIEFS.md, chapter 13.
\chapter{Asymptotic Analysis and Recurrences}\label{ch:b1:asymptotic-analysis}

Two programs sort thirty million records, and both are correct. On a machine that performs a billion \omterm{def:b1:asymptotic-analysis:ram}{elementary
operations} a second, the first finishes in under a second and the second runs for more than ten days. Neither the
language nor the machine makes the difference: the first does about $n \log_2 n$ operations and the second about
$n^2$. This chapter makes such statements precise: what is counted, how \omterm{def:b1:asymptotic-analysis:cost}{costs} are compared when $n$ grows, how the
\omterm{def:b1:asymptotic-analysis:cost}{cost} of an algorithm differs between its best, worst and average inputs, and how to solve the recurrences that the
\omterm{def:b1:asymptotic-analysis:cost}{cost} of a recursive algorithm obeys.

\begin{recall}
For sequences $u$ and $v$ of positive reals, $\lim u_n/v_n = 0$ means that $u$ becomes negligible beside $v$. For
$a > 0$ and $b > 1$: $(\log n)/n^a \to 0$, $n^a/b^n \to 0$ and $b^n/n! \to 0$. The sum $1 + 2 + \dots + n = n(n + 1)/2$,
and $1 + q + \dots + q^{m - 1} = (q^m - 1)/(q - 1)$ for $q \neq 1$. Insertion sort inserts each element into the sorted
part before it, shifting the larger elements one place to the right.
\end{recall}

\section{What the model counts}\label{sec:b1:asymptotic-analysis:model}

\begin{definition}[Model of computation, elementary operation]\label{def:b1:asymptotic-analysis:ram}
The \emph{RAM model}\index{RAM model} is a model of computation: an idealised machine with a memory of numbered
cells, each holding a machine word, in which each \emph{elementary operation}\index{elementary operation} \omterm{def:b1:asymptotic-analysis:cost}{costs} one
unit of time --- an arithmetic operation or a comparison on words, reading or writing a cell, a jump. Words have
$\Theta(\log n)$ bits for inputs of size $n$, enough to index the input.
\end{definition}

\begin{definition}[Input size, cost function]\label{def:b1:asymptotic-analysis:cost}
The \emph{input size}\index{input size} $n$ is a natural number measuring an input: the length of a list, the number of
vertices and edges of a graph, the number of bits of an integer. The \emph{cost}\index{cost function} of an algorithm
on an input is the number of \omterm{def:b1:asymptotic-analysis:ram}{elementary operations} it performs; a \emph{cost function} gives a cost for each size.
\end{definition}

Counting every operation is tedious and depends on details of the translation to machine code, so in practice one
counts a single operation that dominates: the comparisons of a sort, the multiplications of an arithmetic algorithm.
The other operations are at most a constant multiple of it. A counter added to the program makes the count observable.

\omcode{code/bachelor-1/13-asymptotic-analysis/listings/count_ops.py}{6}{18}{Insertion sort, counting its
comparisons}\label{lst:b1:asymptotic-analysis:count-ops}

\omoutput{code/bachelor-1/13-asymptotic-analysis/outputs/count_ops.txt}{Comparisons on a sorted, a reversed and a
shuffled list.}\label{out:b1:asymptotic-analysis:count-ops}

On a sorted list, each insertion stops at its first comparison: $n - 1$ comparisons. On a reversed list, the $i$-th
insertion goes to the front, with $i$ comparisons: $1 + 2 + \dots + (n - 1) = n(n - 1)/2$. A shuffled list lands near
$n(n - 1)/4$, which \cref{sec:b1:asymptotic-analysis:cases} explains. Constants are dropped for two reasons: they depend
on the machine and the \omterm{def:b1:compilation-pipeline:compiler}{compiler}, which the model abstracts away, and they matter less than the growth as soon as $n$
is large. The model also hides the memory hierarchy: on a real machine, an access to a cell far from the last ones
\omterm{def:b1:asymptotic-analysis:cost}{costs} much more than a comparison, which the second year studies.

\section{\texorpdfstring{$O$, $\Omega$ and $\Theta$}{O, Omega and Theta}}\label{sec:b1:asymptotic-analysis:big-o}

\begin{definition}[Big O, big Omega, big Theta]\label{def:b1:asymptotic-analysis:big-o}
Let $f, g : \mathbb{N} \to \mathbb{R}_+$. Then $f \in O(g)$\index{big O} if there are a constant $c > 0$ and a rank $n_0$
with $f(n) \leq c\,g(n)$ for all $n \geq n_0$; $f \in \Omega(g)$\index{big Omega} if $g \in O(f)$, that is $f(n) \geq
c\,g(n)$ from some rank on for some $c > 0$; and $f \in \Theta(g)$\index{big Theta} if both hold.
\end{definition}

One often writes $f = O(g)$; the equality sign is then read as ``belongs to''. The constant $c$ and the rank $n_0$ are
chosen once, before $n$: \cref{exo:b1:asymptotic-analysis:5} shows what goes wrong otherwise. If $f(n)/g(n)$ has a
limit $\ell$, then $\ell < \infty$ gives $f \in O(g)$, and $0 < \ell < \infty$ gives $f \in \Theta(g)$.

\begin{proposition}[Rules for O]\label{prop:b1:asymptotic-analysis:rules}
For \omterm{def:b1:asymptotic-analysis:cost}{cost functions} $f_1 \in O(g_1)$ and $f_2 \in O(g_2)$: $f_1 + f_2 \in O(\max(g_1, g_2))$ and $f_1 f_2 \in O(g_1
g_2)$; if moreover $g_1 \in O(h)$, then $f_1 \in O(h)$. The same rules hold for $\Omega$ and $\Theta$.
\end{proposition}

\begin{proof}
With $f_1 \leq c_1 g_1$ from $n_1$ on and $f_2 \leq c_2 g_2$ from $n_2$ on: from $\max(n_1, n_2)$ on, $f_1 + f_2 \leq
(c_1 + c_2) \max(g_1, g_2)$ and $f_1 f_2 \leq c_1 c_2\, g_1 g_2$. For transitivity, $g_1 \leq c\,h$ from some rank on
gives $f_1 \leq c_1 c\, h$.
\end{proof}

The sum rule says that a sequence of steps \omterm{def:b1:asymptotic-analysis:cost}{costs} as much as its most expensive step, and the product rule that a
loop \omterm{def:b1:asymptotic-analysis:cost}{costs} its number of iterations times the \omterm{def:b1:asymptotic-analysis:cost}{cost} of one iteration, when that \omterm{def:b1:asymptotic-analysis:cost}{cost} does not depend on the iteration.

\begin{proposition}[The scale of growth]\label{prop:b1:asymptotic-analysis:scale}
For $a > 0$ and $b > 1$, as $n \to \infty$: $\log n \in O(n^a)$, $n^a \in O(b^n)$ and $b^n \in O(n!)$, and none of these
is a $\Theta$. Moreover $\log_b n \in \Theta(\log_2 n)$, so the base of a logarithm does not matter inside $O$.
\end{proposition}

\begin{proof}
The three limits recalled at the start of the chapter are 0, which gives each $O$ and excludes each $\Theta$. And
$\log_b n = \log_2 n / \log_2 b$ is a constant multiple.
\end{proof}

\omoutput{code/bachelor-1/13-asymptotic-analysis/outputs/growth.txt}{The usual \omterm{def:b1:asymptotic-analysis:cost}{cost functions}, and the two programs of
the opening at thirty million records and a billion operations a
second.}\label{out:b1:asymptotic-analysis:growth}

\begin{omfigure}
\begin{tikzpicture}
\begin{axis}[width=11cm, height=6.4cm, ymode=log, xmin=1, xmax=40, ymin=1, ymax=1e12, xlabel={$n$},
  ylabel={cost (log scale)}, legend pos=outer north east, legend style={font=\small}, tick label style={font=\small},
  domain=1:40, samples=79]
\addplot[thick, omCode] {ln(x)/ln(2) + 1}; \addlegendentry{$1 + \log_2 n$}
\addplot[thick, omDef] {x}; \addlegendentry{$n$}
\addplot[thick, omThm] {x*ln(x)/ln(2) + 1}; \addlegendentry{$1 + n \log_2 n$}
\addplot[thick, dashed] {x^2}; \addlegendentry{$n^2$}
\addplot[thick, dotted] {x^3}; \addlegendentry{$n^3$}
\addplot[very thick, omThm!50, domain=1:39.8] {2^x}; \addlegendentry{$2^n$}
\end{axis}
\end{tikzpicture}
\par\smallskip{\small The usual \omterm{def:b1:asymptotic-analysis:cost}{cost functions} on a logarithmic scale. Powers of $n$ become straight lines of
increasing slope on a log--log plot; here, with $n$ on a linear axis, the exponential $2^n$ is the straight line, and it
crosses every power of $n$ eventually.}
\end{omfigure}

$O$ is only an upper bound. ``Merge sort is in $O(n^3)$'' is true and useless; ``insertion sort is in $O(n^2)$'' holds for
every input, including the sorted ones that it finishes in $n - 1$ comparisons. A lower bound needs $\Omega$, and a
precise statement needs $\Theta$, with the \omterm{def:b1:c-modules-make-tests:input-class}{input class} it is about.

\section{Worst, best and average case}\label{sec:b1:asymptotic-analysis:cases}

\begin{definition}[Worst-case, best-case and average-case cost]\label{def:b1:asymptotic-analysis:cases}
For an algorithm whose \omterm{def:b1:asymptotic-analysis:cost}{cost} on an input $x$ is $C(x)$, and $I_n$ the finite set of inputs of size $n$ considered, the
\emph{worst-case cost}\index{worst-case cost} is $\max_{x \in I_n} C(x)$, the \emph{best-case cost}\index{best-case cost}
is $\min_{x \in I_n} C(x)$, and the \emph{average-case cost}\index{average-case cost} is $\frac{1}{|I_n|}\sum_{x \in I_n}
C(x)$.
\end{definition}

For a comparison sort on distinct keys, only the order of the keys matters, so $I_n$ can be taken as the $n!$ orders of
$0, \dots, n - 1$. For insertion sort, the best case is $n - 1$ and the worst $n(n - 1)/2$; the average needs a count.

\begin{proposition}[Average cost of insertion sort]\label{prop:b1:asymptotic-analysis:insertion}
Over the $n!$ orders of $n$ distinct keys, insertion sort makes on average exactly $\frac{n(n - 1)}{4} + n - H_n$
comparisons, where $H_n = 1 + \frac12 + \dots + \frac1n$. This is in $\Theta(n^2)$.
\end{proposition}

\begin{proof}
An \emph{inversion} of an order is a pair of positions $i < j$ whose keys are in decreasing order. In insertion sort,
each comparison either shifts an element, which removes exactly one inversion, or stops the inner loop. The inner loop
stops on a comparison unless the key reaches the front, which happens exactly when the key is smaller than every key
before it, a \emph{left-to-right minimum}. So an order with $I$ inversions and $L$ left-to-right minima (the first
key counts) \omterm{def:b1:asymptotic-analysis:cost}{costs} $C = I + (n - 1) - (L - 1)$.

Average of $I$: exchanging the keys at two positions $i < j$ is a one-to-one correspondence between the orders where
the pair is inverted and those where it is not, so each of the $\binom{n}{2}$ pairs is inverted in exactly half of
the orders: the total of $I$ is $n! \binom{n}{2}/2$, and the average $n(n - 1)/4$. Average of $L$: position $i$ (from 0)
holds a left-to-right minimum when it holds the smallest of the first $i + 1$ keys; grouping the orders by the set of
those keys and the arrangement of the rest, that smallest key is at each of the $i + 1$ positions equally often, so in
exactly $1/(i + 1)$ of the orders. Summing over $i$, the average of $L$ is $H_n$, and the average \omterm{def:b1:asymptotic-analysis:cost}{cost} is $n(n - 1)/4 +
(n - 1) - (H_n - 1)$. As $H_n \leq n$, it lies between $n(n - 1)/4$ and $n(n + 3)/4$.
\end{proof}

\omoutput{code/bachelor-1/13-asymptotic-analysis/outputs/inversions.txt}{The average over all orders, as an exact
fraction, against the formula, for $n \leq 8$.}\label{out:b1:asymptotic-analysis:inversions}

\begin{definition}[Space cost]\label{def:b1:asymptotic-analysis:space}
The \emph{space cost}\index{space cost} of an algorithm is the number of memory cells it uses beyond its input, at its
peak, as a function of the \omterm{def:b1:asymptotic-analysis:cost}{input size}; the \omterm{def:b1:c-functions-and-stack:call-stack}{stack frames} of recursive calls count.
\end{definition}

Insertion sort uses $\Theta(1)$ cells beyond the list: it sorts \emph{in place}. A merge sort that copies its halves
uses $\Theta(n)$ cells, and the recursion of a binary search uses $\Theta(\log n)$ frames where the loop of
\cref{ch:b1:program-correctness} uses $\Theta(1)$.

\section{Sums}\label{sec:b1:asymptotic-analysis:sums}

Nested loops give sums. A double loop over $0 \leq i < j < n$ performs $\binom{n}{2} = n(n - 1)/2$ iterations; a triple
loop over $i < j < k$ performs $\binom{n}{3}$ (\cref{exo:b1:asymptotic-analysis:7}). A loop that doubles $j$ from 1
until $n$ runs $\lceil \log_2 n \rceil$ times. Three tools bound the sums that have no closed form.

\begin{method}[Bounding a sum]\label{met:b1:asymptotic-analysis:sums}
\begin{enumerate}
\item \emph{Bound each term by the largest}: $\sum_{i=1}^{n} f(i) \leq n \max_i f(i)$, an upper bound.
\item \emph{Split}: for $f$ increasing, the $\lceil n/2 \rceil$ terms from $i = \lfloor n/2 \rfloor + 1$ on are at least
$f(n/2)$, so $\sum_{i=1}^n f(i) \geq \frac{n}{2} f(\frac{n}{2})$, a lower bound of the same order as the upper one when
$f(n/2) \in \Theta(f(n))$.
\item \emph{Compare with an integral}: for $f$ decreasing on $[1, \infty)$, $\int_1^{n + 1} f \leq \sum_{i=1}^n f(i)
\leq f(1) + \int_1^n f$.
\end{enumerate}
\end{method}

\begin{proposition}[Harmonic numbers and $\log n!$]\label{prop:b1:asymptotic-analysis:harmonic}
$\ln(n + 1) \leq H_n \leq 1 + \ln n$, so $H_n \in \Theta(\log n)$; and $\log_2 n! \in \Theta(n \log n)$.
\end{proposition}

\begin{proof}
The integral of $1/x$ from 1 to $n + 1$ is $\ln(n + 1)$, and from 1 to $n$ it is $\ln n$; the third rule gives the first
claim. For the second, $\log_2 n! = \sum_{i=1}^n \log_2 i \leq n \log_2 n$ by the first rule, and $\geq \frac{n}{2}
\log_2 \frac{n}{2} = \frac{n}{2}(\log_2 n - 1)$ by the second, which is at least $\frac{n}{4}\log_2 n$ for $n \geq 4$.
\end{proof}

The harmonic number thus contributes only $\Theta(\log n)$ to the average of insertion sort, and $\log_2 n!$ will be the
number of comparisons that any comparison sort needs in the worst case (\cref{ch:b1:sorting}).

\section{Recurrences}\label{sec:b1:asymptotic-analysis:recurrences}

\begin{definition}[Recurrence relation, recursion tree]\label{def:b1:asymptotic-analysis:recurrence}
A \emph{recurrence relation}\index{recurrence relation} defines a \omterm{def:b1:asymptotic-analysis:cost}{cost function} by its values on small sizes and an
equation expressing $T(n)$ in terms of values on smaller sizes. The \emph{recursion tree}\index{recursion tree} of such
an equation has a node for each call, labelled with the \omterm{def:b1:asymptotic-analysis:cost}{cost} spent in that call outside its recursive calls; its
children are the calls it makes. $T(n)$ is the sum of all the labels, which one often computes level by level.
\end{definition}

A recursive algorithm gives a recurrence directly: binary search halves its interval with one comparison, $T(n) =
T(n/2) + 1$; merge sort sorts two halves and merges them with at most $n$ comparisons, $T(n) = 2T(n/2) + n$. Three
methods solve them. \emph{Unrolling} substitutes the equation into itself: $T(n) = T(n/2) + 1 = T(n/4) + 2 = \dots =
T(1) + \log_2 n$ for $n$ a power of 2. The \emph{\omterm{def:b1:asymptotic-analysis:recurrence}{recursion tree}} sums level by level. And one can \emph{guess and prove
by induction}, the guess coming from either method or from computed values.

\begin{omfigure}
\begin{tikzpicture}[n/.style={draw, rounded corners=2pt, fill=omCode!8, font=\scriptsize, inner sep=2pt,
  minimum width=9mm}, level distance=11mm]
\node[n] (r) at (0,0) {$n$};
\node[n] (a) at (-3,-1.1) {$n/2$}; \node[n] (b) at (3,-1.1) {$n/2$};
\foreach \p/\q in {a/-4.5, a/-1.5, b/1.5, b/4.5} {\draw (\p) -- (\q,-2.2) node[n] {$n/4$};}
\draw (r) -- (a); \draw (r) -- (b);
\foreach \x in {-5.25,-3.75,-2.25,-0.75,0.75,2.25,3.75,5.25} {\node[font=\scriptsize] at (\x,-3.1) {$\vdots$};}
\foreach \x in {-5.4,-4.6,...,5.4} {\node[draw, fill=omDef!20, inner sep=1.3pt, font=\tiny] at (\x,-3.9) {1};}
\node[font=\scriptsize, anchor=west] at (6.2,0) {level 0: $n$};
\node[font=\scriptsize, anchor=west] at (6.2,-1.1) {level 1: $2 \cdot n/2 = n$};
\node[font=\scriptsize, anchor=west] at (6.2,-2.2) {level 2: $4 \cdot n/4 = n$};
\node[font=\scriptsize, anchor=west] at (6.2,-3.9) {level $\log_2 n$: $n$ leaves};
\node[font=\scriptsize, anchor=west, text=omThm] at (6.2,-4.7) {total: $n \log_2 n$ (+ leaves)};
\end{tikzpicture}
\par\smallskip{\small The \omterm{def:b1:asymptotic-analysis:recurrence}{recursion tree} of $T(n) = 2T(n/2) + n$: every level \omterm{def:b1:asymptotic-analysis:cost}{costs} $n$, and there are $\log_2 n$
levels above the leaves, hence $T(n) = n\log_2 n + nT(1)$ for $n$ a power of 2.}
\end{omfigure}

\omoutput{code/bachelor-1/13-asymptotic-analysis/outputs/recurrences.txt}{Four recurrences computed exactly for $n$ a
power of 2, each equal to its closed form: $T_1 = \log_2 n$, $T_2 = n\log_2 n$, $T_3 = 3n^{\log_2 3} - 2n$ and $T_4 =
2n^2 - n$.}\label{out:b1:asymptotic-analysis:recurrences}

The forms of $T_3$ and $T_4$ come from the proof below; checking them by induction is short. For $T_4$: $T_4(1) = 1 = 2
- 1$, and $4(2(n/2)^2 - n/2) + n = 2n^2 - 2n + n = 2n^2 - n$.

\begin{theorem}[Master theorem]\label{thm:b1:asymptotic-analysis:master}
Let $a \geq 1$, $b \geq 2$ an integer, $k \geq 0$, and $T$ defined on the powers of $b$ by $T(1) > 0$ and $T(n) =
aT(n/b) + n^k$. Then\index{master theorem}
\[
T(n) \in \begin{cases} \Theta(n^k) & \text{if } a < b^k \quad\text{(the root dominates)},\\
\Theta(n^k \log n) & \text{if } a = b^k \quad\text{(every level costs the same)},\\
\Theta(n^{\log_b a}) & \text{if } a > b^k \quad\text{(the leaves dominate)}. \end{cases}
\]
\end{theorem}

\begin{proof}
Let $n = b^m$. Unrolling, level $i$ of the \omterm{def:b1:asymptotic-analysis:recurrence}{recursion tree} has $a^i$ nodes of size $n/b^i$, each costing $(n/b^i)^k$,
and the $a^m$ leaves \omterm{def:b1:asymptotic-analysis:cost}{cost} $T(1)$ each:
\[
T(n) = a^m T(1) + \sum_{i=0}^{m-1} a^i \Bigl(\frac{n}{b^i}\Bigr)^k = a^m T(1) + n^k \sum_{i=0}^{m-1} q^i, \qquad q =
\frac{a}{b^k}.
\]
Note that $a^m = b^{m \log_b a} = n^{\log_b a}$. If $a < b^k$, then $q < 1$ and $1 \leq \sum q^i \leq 1/(1 - q)$; and
$\log_b a < k$, so $a^m T(1) \leq T(1) n^k$; hence $n^k \leq T(n) \leq (T(1) + 1/(1 - q))\, n^k$. If $a > b^k$, then
$q > 1$ and $n^k \sum_{i<m} q^i = n^k (q^m - 1)/(q - 1) = (a^m - n^k)/(q - 1)$, between 0 and $a^m/(q - 1)$; so $T(1)\,
a^m \leq T(n) \leq (T(1) + 1/(q - 1))\, a^m$. The case $a = b^k$ is \cref{exo:b1:asymptotic-analysis:10}.
\end{proof}

\begin{remark}
Algorithms split $n$ into $\lfloor n/b \rfloor$ and $\lceil n/b \rceil$ parts, and their \omterm{def:b1:asymptotic-analysis:cost}{costs} are $an^k$ plus lower
terms. The conclusions still hold for such recurrences when $T$ is increasing: $T(n)$ lies between its values at the
powers of $b$ around $n$, which are within a constant factor of each other. \Cref{exo:b1:asymptotic-analysis:11} treats
merge sort's recurrence with floors and ceilings exactly.
\end{remark}

\omcode{code/bachelor-1/13-asymptotic-analysis/listings/master.py}{6}{14}{The three
cases}\label{lst:b1:asymptotic-analysis:master}

\omoutput{code/bachelor-1/13-asymptotic-analysis/outputs/master.txt}{The cases on six recurrences, and the ratio of the
exact $T(n)$, with $T(1) = 1$, to the predicted order at $2^{10}$ and $2^{20}$: it settles to a
constant.}\label{out:b1:asymptotic-analysis:master}

Karatsuba's multiplication of $n$-digit numbers replaces four half-size products by three, which moves the \omterm{def:b1:asymptotic-analysis:cost}{cost} from
$\Theta(n^2)$ to $\Theta(n^{\log_2 3}) \approx \Theta(n^{1.585})$; Strassen's product of matrices replaces eight
half-size products by seven, from $\Theta(n^3)$ to $\Theta(n^{2.807})$. Both are leaf-dominated: the count of
subproblems matters more than the \omterm{def:b1:asymptotic-analysis:cost}{cost} of combining them. The ratio for merge sort, $1 + 1/\log_2 n$, shows how slowly
the lower term fades.

\begin{omfigure}
\begin{tikzpicture}[x=1cm, y=1cm]
\foreach \col/\title/\a/\b/\c/\note in {0/{$a < b^k$: $2T(n/2) + n^2$}/{$n^2$}/{$n^2/2$}/{$n^2/4$}/{levels shrink:
root},
  5.2/{$a = b^k$: $2T(n/2) + n$}/{$n$}/{$n$}/{$n$}/{levels equal: $\times \log n$},
  10.4/{$a > b^k$: $4T(n/2) + n$}/{$n$}/{$2n$}/{$4n$}/{levels grow: leaves}} {
  \node[font=\scriptsize\bfseries] at (\col+1.6,0.7) {\title};
  \node[font=\scriptsize, anchor=west] at (\col+2.9,0.05) {\a};
}
\fill[omCode!25] (0.4,-0.1) rectangle (2.8,0.2); \fill[omCode!25] (5.6,-0.1) rectangle (8.0,0.2);
\fill[omCode!25] (0.4,-0.6) rectangle (1.6,-0.3); \node[font=\scriptsize, anchor=west] at (2.9,-0.45) {$n^2/2$};
\fill[omCode!25] (0.4,-1.1) rectangle (1.0,-0.8); \node[font=\scriptsize, anchor=west] at (2.9,-0.95) {$n^2/4$};
\fill[omCode!25] (0.4,-1.6) rectangle (0.7,-1.3); \node[font=\scriptsize, anchor=west] at (2.9,-1.45) {$\cdots$};
\fill[omCode!25] (5.6,-0.6) rectangle (8.0,-0.3); \node[font=\scriptsize, anchor=west] at (8.1,-0.45) {$n$};
\fill[omCode!25] (5.6,-1.1) rectangle (8.0,-0.8); \node[font=\scriptsize, anchor=west] at (8.1,-0.95) {$n$};
\fill[omCode!25] (5.6,-1.6) rectangle (8.0,-1.3); \node[font=\scriptsize, anchor=west] at (8.1,-1.45) {$\cdots$};
\fill[omThm!30] (10.8,-0.6) rectangle (11.4,-0.3); \node[font=\scriptsize, anchor=west] at (13.3,-0.45) {$2n$};
\fill[omThm!30] (10.8,-1.1) rectangle (12.0,-0.8); \node[font=\scriptsize, anchor=west] at (13.3,-0.95) {$4n$};
\fill[omThm!30] (10.8,-1.6) rectangle (13.2,-1.3); \node[font=\scriptsize, anchor=west] at (13.3,-1.45) {$\cdots
n^2$};
\fill[omThm!30] (10.8,-0.1) rectangle (11.1,0.2);
\node[font=\scriptsize, text=omThm] at (2,-2.1) {total $\Theta(n^2)$};
\node[font=\scriptsize, text=omThm] at (7.2,-2.1) {total $\Theta(n \log n)$};
\node[font=\scriptsize, text=omThm] at (12.2,-2.1) {total $\Theta(n^2)$};
\end{tikzpicture}
\par\smallskip{\small The three regimes of the master theorem, drawn as the \omterm{def:b1:asymptotic-analysis:cost}{cost} of each level of the \omterm{def:b1:asymptotic-analysis:recurrence}{recursion tree}
(one bar per level, from the root down). The level \omterm{def:b1:asymptotic-analysis:cost}{costs} form a geometric sequence of ratio $a/b^k$: decreasing, the
root's term dominates; constant, the $\log_b n$ levels add up; increasing, the leaves dominate.}
\end{omfigure}

\section{Counts and times}\label{sec:b1:asymptotic-analysis:times}

A count predicts a time if each counted operation \omterm{def:b1:asymptotic-analysis:cost}{costs} about the same. The \emph{doubling test} checks the order of a
\omterm{def:b1:asymptotic-analysis:cost}{cost} from measurements alone: double $n$ and look at the ratio of the \omterm{def:b1:asymptotic-analysis:cost}{costs}. For a \omterm{def:b1:asymptotic-analysis:cost}{cost} in $\Theta(n^2)$ the ratio tends
to 4, for $\Theta(n)$ to 2, and for $\Theta(n \log n)$ to $2(1 + 1/\log_2 n)$, slightly above 2 (the weekend problem
proves it).

\omcode{code/bachelor-1/13-asymptotic-analysis/listings/doubling.py}{8}{20}{Counts of three algorithms as $n$
doubles}\label{lst:b1:asymptotic-analysis:doubling}

\omoutput{code/bachelor-1/13-asymptotic-analysis/outputs/doubling.txt}{Comparisons of the maximum, merge sort and
insertion sort on shuffled lists, then the ratios to the previous row.}\label{out:b1:asymptotic-analysis:doubling}

\omoutput{code/bachelor-1/13-asymptotic-analysis/outputs/timing.txt}{The same three algorithms timed on the author's
laptop, best of five runs, in milliseconds, then the ratios.}\label{out:b1:asymptotic-analysis:timing}

\begin{omfigure}
\begin{tikzpicture}
\begin{axis}[width=10cm, height=5.4cm, xmode=log, log basis x=2, xmin=400, xmax=10000, ymin=0, ymax=130,
  xtick={500,1000,2000,4000,8000}, xticklabels={500,1000,2000,4000,8000}, xlabel={$n$},
  ylabel={ns per comparison}, legend pos=outer north east, legend style={font=\small},
  tick label style={font=\small}]
\addplot[thick, mark=*, omThm] table[x=n, y=merge]
  {code/bachelor-1/13-asymptotic-analysis/figures/timing.dat}; \addlegendentry{merge sort}
\addplot[thick, mark=square*, omDef] table[x=n, y=max]
  {code/bachelor-1/13-asymptotic-analysis/figures/timing.dat}; \addlegendentry{maximum}
\addplot[thick, mark=triangle*, omCode] table[x=n, y=ins]
  {code/bachelor-1/13-asymptotic-analysis/figures/timing.dat}; \addlegendentry{insertion sort}
\end{axis}
\end{tikzpicture}
\par\smallskip{\small Time divided by count, from the two outputs above: nearly flat for each algorithm, so the count
predicts the time up to a constant, which differs between algorithms (a merge step allocates and appends, an insertion
step only compares and moves).}
\end{omfigure}

The time ratios scatter around the count ratios: a measurement on a real machine includes caches, the \omterm{def:b1:compilation-pipeline:interpreter}{interpreter}
and the other programs running at the same time. The counts are exact and repeatable; the times say what the counts
\omterm{def:b1:asymptotic-analysis:cost}{cost} on one machine.

\begin{inpractice}
To estimate an unknown \omterm{def:b1:asymptotic-analysis:cost}{cost}, run the program for sizes $n, 2n, 4n, \dots$ until the time is a second or more, and look
at the ratios. A ratio near 2 suggests linear time; near 4 quadratic; near 8 cubic. A ratio that keeps growing suggests
exponential time, and the end of the experiment.
\end{inpractice}

\section{Exercises}\label{sec:b1:asymptotic-analysis:exercises}

\begin{exercise}[$\star$]\label{exo:b1:asymptotic-analysis:1}
Predict what the program prints, and give the two counts as functions of the bounds 10 and 1000.
\omcode{code/bachelor-1/13-asymptotic-analysis/listings/exo_count.py}{3}{12}{Two counts}
\end{exercise}

\begin{exercise}[$\star$]\label{exo:b1:asymptotic-analysis:2}
Order by growth, from slowest to fastest: $n^2 \log n$, $(\log n)^{10}$, $1.01^n$, $\sqrt{n}$, $n^{2.5}$. Justify each
step with \cref{prop:b1:asymptotic-analysis:scale}.
\end{exercise}

\begin{exercise}[$\star$]\label{exo:b1:asymptotic-analysis:3}
With $T(1) = 1$ and $T(n) = 2T(n/2) + 1$, compute $T(16)$, guess a closed form, and prove it.
\end{exercise}

\begin{exercise}[$\star$]\label{exo:b1:asymptotic-analysis:4}
Which case of the master theorem applies, and what is the order of $T$? (a) $T(n) = 4T(n/2) + n$. (b) $T(n) = 4T(n/2)
+ n^2$. (c) $T(n) = T(n/3) + n$. (d) $T(n) = 2T(n/4) + \sqrt{n}$.
\end{exercise}

\begin{exercise}[$\star\star$]\label{exo:b1:asymptotic-analysis:5}
Find the error. ``We prove $n^2 \in O(n)$ by induction on $n$. For $n = 1$, $1 \leq c \cdot 1$ with $c = 1$. If $n^2
\leq c n$, then $(n + 1)^2 = n^2 + 2n + 1 \leq cn + 3n = (c + 3)n \leq c'(n + 1)$ with $c' = c + 3$. So $n^2 \in O(n)$.''
\end{exercise}

\begin{exercise}[$\star\star$]\label{exo:b1:asymptotic-analysis:6}
A student reasons: ``the inner loop runs at most $\log_2 n$ times, so the function \omterm{def:b1:asymptotic-analysis:cost}{costs} $\Theta(n \log n)$''. What is
right and what is wrong? Find the true order, and prove it.
\omcode{code/bachelor-1/13-asymptotic-analysis/listings/exo_cost.py}{6}{13}{A loop that depends on $i$}
\omoutput[2-4]{code/bachelor-1/13-asymptotic-analysis/outputs/exo_cost.txt}{The counts.}
\end{exercise}

\begin{exercise}[$\star\star$]\label{exo:b1:asymptotic-analysis:7}
Write \code{count_triples(n)}, which counts the iterations of three nested loops over $0 \leq i < j < k < n$. Find a
closed form, prove it, and test it against the counter.
\end{exercise}

\begin{exercise}[$\star\star$]\label{exo:b1:asymptotic-analysis:8}
Prove that $\sum_{i=1}^n i \log_2 i \in \Theta(n^2 \log n)$.
\end{exercise}

\begin{exercise}[$\star\star$]\label{exo:b1:asymptotic-analysis:9}
Solve exactly $T(0) = 0$, $T(n) = T(n - 1) + n$, and $U(0) = 0$, $U(n) = 2U(n - 1) + 1$. Why does the master theorem not
apply?
\end{exercise}

\begin{exercise}[$\star\star\star$]\label{exo:b1:asymptotic-analysis:10}
Prove the case $a = b^k$ of \cref{thm:b1:asymptotic-analysis:master}: give $T(n)$ exactly for $n = b^m$, and deduce
$T(n) \in \Theta(n^k \log n)$.
\end{exercise}

\begin{exercise}[$\star\star\star$]\label{exo:b1:asymptotic-analysis:11}
Let $T(1) = 0$ and $T(n) = T(\lfloor n/2 \rfloor) + T(\lceil n/2 \rceil) + n$ for $n \geq 2$. Prove that $T$ is increasing
and that $T(n) \in \Theta(n \log n)$; then check that $T(n) = n\lceil \log_2 n \rceil + n - 2^{\lceil \log_2 n
\rceil}$ for every $n \geq 1$.
\end{exercise}

\begin{exercise}[$\star\star\star$]\label{exo:b1:asymptotic-analysis:12}
Write \code{identify(a, b, k)}, which computes $T(b^i)$ for $T(n) = aT(n/b) + n^k$ and $T(1) = 1$, estimates the
exponent of $T$ from the numbers, and says which case of the master theorem the numbers show, without using the
comparison of $a$ with $b^k$. Test it against \cref{lst:b1:asymptotic-analysis:master}.
\end{exercise}

\section{Problem: Three ways to find a duplicate}\label{sec:b1:asymptotic-analysis:problem}

\begin{problem}[{Weekend problem --- Three ways to find a duplicate}]\label{pb:b1:asymptotic-analysis:1}
The town's library keeps its loans in order; which books were borrowed more than once among the first $n$ loans?
Three functions answer, each counting its operations: one comparison of two loans, one comparison of neighbours after
sorting, or one look-up in a set.
\omcodes{code/bachelor-1/13-asymptotic-analysis/listings/duplicates.py}{14-42}{Three ways to find the
duplicates}\label{lst:b1:asymptotic-analysis:duplicates}
\omoutput{code/bachelor-1/13-asymptotic-analysis/outputs/duplicates.txt}{The number of duplicated books, the
counts, and their ratios as $n$ doubles.}\label{out:b1:asymptotic-analysis:duplicates}

\medskip\noindent\textbf{Part I --- The double loop.}
\begin{enumerate}
\item Show that \code{by_pairs} makes exactly $n(n - 1)/2$ comparisons, whatever the data.
\item Deduce that it is in $\Theta(n^2)$, in the best, the worst and the average case.
\item Check the column \code{pairs} of the table.
\item Compute the ratio of the counts for $2n$ and $n$, and its limit.
\item Could the function stop early? Rewrite its specification so that it can (``is there a duplicate?''), and give its
best and worst cases then.
\item What is its \omterm{def:b1:asymptotic-analysis:space}{space cost}, beyond the input and the answer?
\end{enumerate}

\medskip\noindent\textbf{Part II --- Sort, then scan.}
\begin{enumerate}[resume]
\item Show that after sorting, two equal values are neighbours, and that the scan finds every duplicated value.
\item Merge sort makes at most $T(n)$ comparisons, with $T(1) = 0$ and $T(n) = T(\lfloor n/2 \rfloor) + T(\lceil n/2
\rceil) + n - 1$. Show that $T(n) \leq n \lceil \log_2 n \rceil$.
\item Deduce that \code{by_sorting} is in $O(n \log n)$.
\item On the table, compare the counts with $n \log_2 n$ for $n = 400$.
\item Show that $(2n)\log_2(2n) / (n \log_2 n) = 2(1 + 1/\log_2 n)$, and compute it for $n = 25$ and $n = 200$.
\item Why are the measured ratios not exactly these values?
\end{enumerate}

\medskip\noindent\textbf{Part III --- A set of seen values.}
\begin{enumerate}[resume]
\item Show that \code{by_set} makes exactly $n$ look-ups.
\item Under which assumption on the set is its \omterm{def:b1:asymptotic-analysis:cost}{cost} in $\Theta(n)$? (Hash tables are studied in
\cref{ch:b1:hash-tables}.)
\item What is the \omterm{def:b1:asymptotic-analysis:space}{space cost} of each of the three methods?
\item The book identifiers are numbers from 1 to 80. Give a method that uses a list of 80 booleans instead of a set, and
its \omterm{def:b1:asymptotic-analysis:cost}{cost}.
\item For $n = 400$, how many times fewer operations does the set make than the double loop?
\item Why does the table count only one operation per look-up, when a look-up does several \omterm{def:b1:asymptotic-analysis:ram}{elementary operations}?
\end{enumerate}

\medskip\noindent\textbf{Part IV --- The doubling ratio.}
\begin{enumerate}[resume]
\item For a \omterm{def:b1:asymptotic-analysis:cost}{cost} $f(n) = c\,n^k$, show that $f(2n)/f(n) = 2^k$ exactly.
\item For $f(n) = c\, n^k + d\,n^{k - 1}$ with $c > 0$, show that the ratio tends to $2^k$.
\item For $f \in \Theta(g)$ only, can the ratio of $f$ fail to converge? Give an example.
\item Read the ratios of the table, and say which order each column suggests.
\item Why do the ratios of the merge sort column approach 2 so slowly?
\item State the named result of the problem, \emph{the doubling ratio test}: the limits of the ratios for counts
$c\,n^2$, $c\,n \log_2 n$ and $c\,n$, proved, and the values read from the table.
\end{enumerate}
\end{problem}
